\chapter{Допуск связного трёхчастичного барионного ядра}
\label{ch:v8-baryon-connected-three-body-kernel-admission-gate}

Предыдущие главы обнаружили отдельно свободу межкопийной ковариации,
пространственного ядра, магнитного контакта и перестановочной ветви. Проверим,
являются ли они частными проявлениями одного пропущенного корреляционного
объекта.

\section{Точная потеря третьей кумулянты}

Пусть $x_1,x_2,x_3\in\{-1,1\}$. На восьми точках введём семейство
\begin{equation}
 p_\theta(x_1,x_2,x_3)
 =\frac18\bigl(1+\theta x_1x_2x_3\bigr),
 \qquad -1\leq\theta\leq1.
 \label{eq:v8-baryon-three-body-parity-family}
\end{equation}
Оно неотрицательно и нормировано. Для всех $\theta$ одно- и двухточечные
моменты совпадают:
\begin{equation}
 \mathbb E_\theta[x_i]=0,
 \qquad
 \mathbb E_\theta[x_ix_j]=0\quad(i\ne j),
 \label{eq:v8-baryon-three-body-lower-moments}
\end{equation}
тогда как связный трёхточечный момент равен
\begin{equation}
 \mathbb E_\theta[x_1x_2x_3]=\theta.
 \label{eq:v8-baryon-three-body-connected-moment}
\end{equation}

Более сильная формулировка использует отображение всех трёх попарных
маргиналов. Его матрица имеет размер $12\times8$, ранг $7$ и одномерное
ядро
\begin{equation}
 \ker M_{\leq2}
 =\operatorname{span}\{(x_1x_2x_3)_{(x_1,x_2,x_3)}\}.
 \label{eq:v8-baryon-pair-marginal-kernel}
\end{equation}
Следовательно, знание всех однокопийных и попарных ограничений не определяет
полное трёхчастичное состояние.

\begin{theorem}[Неинъективность двухточечного барионного чтения]
Отображение из полного трёхчастичного состояния в совокупность всех его
одно- и двухчастичных ограничений неинъективно уже на диагональной
восьмиточечной алгебре. Поэтому полный двухточечный процесс Тома VIII не
может сам по себе определить связное трёхчастичное ядро.
\end{theorem}

\section{Что сохраняет слабый столкновительный предел}

Полный микроскопический контакт предыдущих глав имеет звёздный вид
\begin{equation}
 H_{\rm int}=\sum_{a=1}^{42}F_a\otimes
 \bigl(|a\rangle\langle0|+|0\rangle\langle a|\bigr).
 \label{eq:v8-baryon-star-interaction-recall}
\end{equation}
Для чётности среды
$\Pi=|0\rangle\langle0|-\sum_a|a\rangle\langle a|$ выполнено
\begin{equation}
 \Pi H_{\rm int}\Pi=-H_{\rm int},
 \qquad
 \langle0|H_{\rm int}^{2m+1}|0\rangle=0.
 \label{eq:v8-baryon-star-odd-moment-zero}
\end{equation}
Таким образом, текущий родитель не содержит независимой третьей вакуумной
кумулянты.

Даже если заменить среду на состояние с ограниченной третьей кумулянтой,
обычный слабый предел её стирает. При $h=\varepsilon^2$ и
$n=u/\varepsilon^2$ накопленные второй и третий порядки масштабируются как
\begin{equation}
 n h=u,
 \qquad
 n h^{3/2}\kappa_3=u\varepsilon\kappa_3\longrightarrow0.
 \label{eq:v8-baryon-third-cumulant-collision-scaling}
\end{equation}
Поэтому гауссовский квантово-марковский предел сохраняет ковариацию, но не
ограниченную связную третью кумулянту.

\section{Операторный уровень продолжения}

В трёхчастичной полевой постановке требуемая амплитуда должна удовлетворять
уравнению вида
\begin{equation}
 \Psi=K_{(3)}\Psi,
 \qquad
 K_{(3)}=K_{(3)}^{\rm irr}+\sum_{a=1}^{3}K_{(2)}^{(a)}.
 \label{eq:v8-baryon-faddeev-kernel-target}
\end{equation}
Эта формула задаёт тип объекта, но не считается его выводом. Текущий
$42$-скачковый процесс определяет второй порядок и не задаёт ни
$K_{(3)}^{\rm irr}$, ни эквивалентную связную трёхчастичную кумулянту.

\begin{nogocriterion}[Запрет продолжать от отдельных коэффициентов]
Параметры межкопийной корреляции, кулоновского матричного элемента,
магнитного контакта и перестановочной ветви нельзя выбирать независимо и
затем называть полным барионным ядром. Ветка открывается только после
родительского вывода связного трёхчастичного оператора или шеститочечного
ядра. Сингулярное усиление третьей кумулянты как
$\kappa_3\sim\varepsilon^{-1}$ также является новым входом, а не следствием
текущего слабого предела.
\end{nogocriterion}

\begin{protocol}
\begin{enumerate}
 \item поле вычислений: \textbf{$\mathbb Q(\theta,\varepsilon,u,\kappa_3)$};
 \item числа с плавающей точкой: \textbf{отсутствуют};
 \item ранг отображения попарных маргиналов: \textbf{$7$ из $8$};
 \item размерность невидимой трёхчастичной свободы: \textbf{$1$};
 \item третья кумулянта семейства: \textbf{$\theta$};
 \item нечётные моменты звёздного родителя: \textbf{нулевые};
 \item ограниченная третья кумулянта переживает предел: \textbf{нет};
 \item двухточечный процесс определяет барионное ядро: \textbf{нет};
 \item требуемый новый объект:
       \textbf{связное трёхчастичное или шеститочечное ядро};
 \item следующий гейт:
       \textbf{происхождение трёхчастичного ядра либо строгий запрет}.
\end{enumerate}
\end{protocol}

Машинный сертификат сохранён в
\path{s2t_v8_baryon_connected_three_body_kernel_admission_gate_results.json}.